\documentclass[a4paper,12pt]{article}
\usepackage[utf8]{inputenc}
\usepackage[russian,english]{babel}
\usepackage[T2A]{fontenc}
\usepackage{amsmath,amssymb}
\usepackage{indentfirst}
\usepackage{array}
\usepackage[pdftex]{graphicx}

\title{Задание 9

С помощью MPI реализовать алгоритм SUMMA перемножения матриц.}
\author{Андрей Преображенский 323 группа}
\date{08.12.2024}

\begin{document}

\maketitle

\newpage

\section{Постановка задачи}
Требуется:
Реализовать параллельный алгоритм SUMMA перемножения матриц для числа процессов, составляющего полный квадрат. Размер сетки N делиться на sqrt(p) * размер блока (B-SIZE). Составить график T(p).

\section{Реализация}
При реализации распараллеливания будем использовать директивы MPI:
\begin{enumerate}
    \item MPI-Comm-split
    \item MPI-Bcast
    \item MPI-Wtime
\end{enumerate}
Размер транслируемого блока будет определяться как N * B-SIZE / sqrt(p).
Запуск программ будем производить на системе Polus.

\section{Результаты}
При фиксированной ширине сетки N = 600 составим график зависимости T(p) - времени работы программы.
B-SIZE возьмём равным 10.
Возьмём p = 1, 4, 9, 16

График T(p):

\includegraphics[width=27em, height=22em]{"time9.png"}


\section{Выводы}

Программа даёт ускорение при увеличении числа процессов. При подобранном значении блока время на пересылку блоков тратиться меньше, чем на возможное вычисление всей матрицы одним процессом. Алгоритм требует больших допущений в плане размера сетки и числа процессов.

\end{document}